\documentclass{article}
\usepackage{amsmath,amssymb,amsthm}
\usepackage{algorithm,algpseudocode}
\usepackage{hyperref}
\usepackage{enumitem}
\newtheorem{theorem}{Theorem}
\newtheorem{lemma}{Lemma}
\newcommand{\prox}{\operatorname{prox}}
\newcommand{\grad}{\nabla}
\newcommand{\R}{\mathbb{R}}
\newcommand{\G}{\mathcal{G}}
\newcommand{\F}{\mathcal{F}}

\begin{document}

\section*{Algorithm Name}
\textbf{Proximal Gradient Method (PGM) for a non‑convex smooth part $f$ and a (possibly) non‑convex nonsmooth part $g$.}

\section*{Mathematical Formulation}
Let  
\[
\F(x)=f(x)+g(x),\qquad x\in\R^{d},
\]
where  

\begin{itemize}
  \item $f:\R^{d}\rightarrow\R$ is differentiable and $L$–smooth, i.e.
        \[
        \|\grad f(x)-\grad f(y)\|\le L\|x-y\|,\qquad\forall x,y\in\R^{d}.
        \]
  \item $g:\R^{d}\rightarrow\R\cup\{+\infty\}$ is proper, lower‑semicontinuous (no convexity required).
\end{itemize}

For a stepsize $\eta>0$ define the \emph{proximal operator}
\[
\prox_{\eta g}(v):=\arg\min_{u\in\R^{d}}\Big\{ g(u)+\frac{1}{2\eta}\|u-v\|^{2}\Big\}.
\]

The iteration reads
\begin{equation}
\boxed{
x_{k+1}= \prox_{\eta g}\bigl(x_{k}-\eta\,\grad f(x_{k})\bigr)},\qquad k=0,1,2,\dots
\label{eq:update}
\end{equation}
Define the \emph{proximal gradient mapping}
\[
\G_{\eta}(x):=\frac{1}{\eta}\Bigl(x-\prox_{\eta g}\bigl(x-\eta \grad f(x)\bigr)\Bigr).
\tag{1}
\]

\section*{Pseudocode}
\begin{algorithm}[H]
\caption{Proximal Gradient Method}
\begin{algorithmic}[1]
\Require Initial point $x_{0}\in\R^{d}$, stepsize $\eta\in(0,1/L]$,
          maximum iterations $K$ (or tolerance $\varepsilon$).
\For{$k=0$ \textbf{to} $K-1$}
    \State Compute gradient $g_{k}\gets \grad f(x_{k})$.
    \State Take a proximal step 
          \[
          x_{k+1}\gets \prox_{\eta g}\bigl(x_{k}-\eta g_{k}\bigr).
          \]
    \If{$\|\G_{\eta}(x_{k})\|\le \varepsilon$}
        \State \textbf{break}
    \EndIf
\EndFor
\State \textbf{return} $x_{k+1}$.
\end{algorithmic}
\end{algorithm}

\section*{Dry Run Example}
Consider the scalar problem  
\[
f(w)=(w-3)^{2},\qquad g\equiv0,
\]
so $\prox_{\eta g}(v)=v$ (identity).  
Take $w_{0}=0$, stepsize $\eta=0.1$ and run three iterations.

\begin{center}
\begin{tabular}{c|c|c|c}
$k$ & $w_{k}$ & $\grad f(w_{k})$ & $w_{k+1}=w_{k}-\eta\grad f(w_{k})$ \\ \hline
0 & $0$ & $2(0-3)=-6$ & $0-0.1(-6)=0.6$ \\
1 & $0.6$ & $2(0.6-3)=-4.8$ & $0.6-0.1(-4.8)=1.08$ \\
2 & $1.08$ & $2(1.08-3)=-3.84$ & $1.08-0.1(-3.84)=1.464$ \\
\end{tabular}
\end{center}

Objective values:
\[
f(0)=9,\; f(0.6)=5.76,\; f(1.08)=3.6864,\; f(1.464)=2.3049,
\]
showing monotone decrease.

\newpage
\section*{Assumptions}
\begin{enumerate}[label=(A\arabic*)]
  \item \label{A1} \textbf{Smoothness of $f$:} $f$ is differentiable and $L$‑smooth:
  \[
  f(y)\le f(x)+\langle\grad f(x),y-x\rangle+\frac{L}{2}\|y-x\|^{2},
  \quad\forall x,y\in\R^{d}.
  \]
  \item \label{A2} \textbf{Bounded below:} $\F$ is bounded from below, i.e.
        \[
        \F_{\inf}:=\inf_{x\in\R^{d}}\F(x)>-\infty .
        \]
  \item \label{A3} \textbf{Stepsize:} Choose $\eta\in\bigl(0,\,\frac{1}{L}\bigr]$.
  \item \label{A4} \textbf{Existence of proximal operator:} For every $v\in\R^{d}$,
        the minimisation defining $\prox_{\eta g}(v)$ admits at least one solution.
\end{enumerate}

\section*{Main Convergence Theorem}
\begin{theorem}[Stationarity rate of PGM]\label{thm:main}
Under Assumptions \ref{A1}–\ref{A4} and for any integer $K\ge1$, the iterates generated by \eqref{eq:update} satisfy
\begin{equation}
\min_{0\le k\le K-1}\bigl\|\G_{\eta}(x_{k})\bigr\|^{2}
\;\le\;
\frac{2\bigl(\F(x_{0})-\F_{\inf}\bigr)}{\eta K}.
\label{eq:rate}
\end{equation}
Consequently,
\[
\lim_{k\to\infty}\bigl\|\G_{\eta}(x_{k})\bigr\|=0,
\]
i.e. every limit point of $\{x_{k}\}$ is a \emph{critical point} of $\F$ (first‑order stationary).
\end{theorem}

\section*{Theoretical Properties}
\begin{itemize}
  \item \textbf{Descent property:} $\F(x_{k+1})\le \F(x_{k})-\frac{\eta}{2}\|\G_{\eta}(x_{k})\|^{2}$, guaranteeing monotone decrease of the objective.
  \item \textbf{Stability w.r.t.\ stepsize:} Any $\eta\le1/L$ yields the same $O(1/K)$ bound; larger $\eta$ may violate Lemma~\ref{lem:descent}.
  \item \textbf{Comparison to Gradient Descent (GD):} When $g\equiv0$, $\G_{\eta}(x)=\grad f(x)$ and \eqref{eq:rate} reduces to the classical $O(1/K)$ rate for smooth non‑convex GD.
  \item \textbf{Robustness to non‑convex $g$:} The proof only uses the \emph{non‑expansiveness} of the proximal operator, not convexity of $g$.
\end{itemize}

\section*{Discussion}
The theorem shows that the proximal gradient mapping—an exact analogue of the gradient in the composite setting—vanishes at the optimal $O(1/K)$ speed, matching the best known rates for first‑order methods on smooth non‑convex problems. The result holds without convexity of $g$, which is crucial for many modern regularisers (e.g., $\ell_{0}$, SCAD). The constant $2/(\eta)$ in \eqref{eq:rate} reveals a linear sensitivity to the stepsize; choosing the largest admissible $\eta=1/L$ yields the smallest bound.

\newpage
\section*{Preliminaries (Definitions and Lemmas)}
\begin{lemma}[Non‑expansiveness of the proximal operator]\label{lem:nonexp}
For any $\eta>0$ and any $u,v\in\R^{d}$,
\[
\bigl\|\prox_{\eta g}(u)-\prox_{\eta g}(v)\bigr\|\le\|u-v\|.
\]
\end{lemma}
\begin{proof}
The proximal operator is the resolvent of the sub‑differential $\partial g$; resolvents of monotone operators are firmly non‑expansive, which implies the claimed inequality. ∎
\end{proof}

\begin{lemma}[Descent Lemma for $L$‑smooth $f$]\label{lem:descent}
For any $x,y\in\R^{d}$,
\[
f(y)\le f(x)+\langle\grad f(x),y-x\rangle+\frac{L}{2}\|y-x\|^{2}.
\]
\end{lemma}
\begin{proof}
Standard consequence of $L$‑Lipschitz continuity of $\grad f$. ∎
\end{proof}

\section*{Main Proof (Step‑by‑Step Derivation)}
We prove Theorem~\ref{thm:main}. Throughout, denote
\[
y_{k}:=\prox_{\eta g}\bigl(x_{k}-\eta\grad f(x_{k})\bigr)=x_{k+1}.
\]

\begin{enumerate}
\item[Step 1.] \textbf{Optimality condition of the proximal step.}  
By definition of $y_{k}$,
\[
y_{k}=\arg\min_{u}\Bigl\{g(u)+\frac{1}{2\eta}\|u-(x_{k}-\eta\grad f(x_{k}))\|^{2}\Bigr\}.
\]
First‑order optimality yields
\begin{equation}
0\in\partial g(y_{k})+\frac{1}{\eta}\bigl(y_{k}-(x_{k}-\eta\grad f(x_{k}))\bigr).
\tag{2}
\end{equation}
Hence there exists $s_{k}\in\partial g(y_{k})$ such that
\[
s_{k}= \frac{1}{\eta}\bigl(x_{k}-\eta\grad f(x_{k})-y_{k}\bigr).
\tag{3}
\]

\item[Step 2.] \textbf{Express the proximal gradient mapping.}  
From \eqref{eq:update} and \eqref{eq:rate} definition,
\[
\G_{\eta}(x_{k})=\frac{1}{\eta}(x_{k}-y_{k})= \grad f(x_{k})+s_{k}.
\tag{4}
\]

\item[Step 3.] \textbf{Apply the descent lemma to $f$ with $x=x_{k}$ and $y=y_{k}$.}
\[
f(y_{k})\le f(x_{k})+\langle\grad f(x_{k}),y_{k}-x_{k}\rangle+\frac{L}{2}\|y_{k}-x_{k}\|^{2}.
\tag{5}
\]

\item[Step 4.] \textbf{Add $g(y_{k})$ to both sides of \eqref{eq:5}.}  
Since $g(y_{k})\le g(y_{k})$ (trivial),
\[
\F(y_{k})\le f(x_{k})+g(y_{k})+\langle\grad f(x_{k}),y_{k}-x_{k}\rangle+\frac{L}{2}\|y_{k}-x_{k}\|^{2}.
\tag{6}
\]

\item[Step 5.] \textbf{Use the sub‑gradient inequality for $g$.}  
From $s_{k}\in\partial g(y_{k})$ we have
\[
g(x_{k})\ge g(y_{k})+\langle s_{k},x_{k}-y_{k}\rangle.
\tag{7}
\]
Rearranging,
\[
g(y_{k})\le g(x_{k})+\langle s_{k},y_{k}-x_{k}\rangle .
\tag{8}
\]

\item[Step 6.] \textbf{Insert \eqref{eq:8} into \eqref{eq:6}.}
\[
\F(y_{k})\le f(x_{k})+g(x_{k})+\langle\grad f(x_{k})+s_{k},y_{k}-x_{k}\rangle+\frac{L}{2}\|y_{k}-x_{k}\|^{2}.
\tag{9}
\]

\item[Step 7.] \textbf{Replace $\grad f(x_{k})+s_{k}$ using \eqref{eq:4}.}
\[
\F(y_{k})\le \F(x_{k})-\langle\G_{\eta}(x_{k}),x_{k}-y_{k}\rangle+\frac{L}{2}\|y_{k}-x_{k}\|^{2}.
\tag{10}
\]

\item[Step 8.] \textbf{Rewrite the inner product.}  
Since $x_{k}-y_{k}= \eta\G_{\eta}(x_{k})$ (from \eqref{eq:4}),
\[
\langle\G_{\eta}(x_{k}),x_{k}-y_{k}\rangle
= \eta\|\G_{\eta}(x_{k})\|^{2}.
\tag{11}
\]

\item[Step 9.] \textbf{Replace $\|y_{k}-x_{k}\|^{2}$ by $\eta^{2}\|\G_{\eta}(x_{k})\|^{2}$.}
\[
\|y_{k}-x_{k}\|^{2}= \eta^{2}\|\G_{\eta}(x_{k})\|^{2}.
\tag{12}
\]

\item[Step 10.] \textbf{Plug \eqref{eq:11} and \eqref{eq:12} into \eqref{eq:10}.}
\[
\F(x_{k+1})\le \F(x_{k})-\eta\|\G_{\eta}(x_{k})\|^{2}
+\frac{L}{2}\eta^{2}\|\G_{\eta}(x_{k})\|^{2}.
\tag{13}
\]

\item[Step 11.] \textbf{Factor the right‑hand side.}
\[
\F(x_{k+1})\le \F(x_{k})-\Bigl(\eta-\frac{L\eta^{2}}{2}\Bigr)\|\G_{\eta}(x_{k})\|^{2}.
\tag{14}
\]

\item[Step 12.] \textbf{Use the stepsize constraint $\eta\le 1/L$.}  
Since $\eta\le 1/L$, we have $L\eta\le1$ and thus
\[
\eta-\frac{L\eta^{2}}{2}\;\ge\;\frac{\eta}{2}.
\tag{15}
\]

\item[Step 13.] \textbf{Obtain the fundamental descent inequality.}
\[
\boxed{\F(x_{k+1})\le \F(x_{k})-\frac{\eta}{2}\|\G_{\eta}(x_{k})\|^{2}}.
\tag{16}
\]

\item[Step 14.] \textbf{Sum \eqref{eq:16} from $k=0$ to $K-1$.}
\[
\F(x_{K})\le \F(x_{0})-\frac{\eta}{2}\sum_{k=0}^{K-1}\|\G_{\eta}(x_{k})\|^{2}.
\tag{17}
\]

\item[Step 15.] \textbf{Use the lower‑boundedness of $\F$.}
Since $\F(x_{K})\ge \F_{\inf}$,
\[
\frac{\eta}{2}\sum_{k=0}^{K-1}\|\G_{\eta}(x_{k})\|^{2}
\le \F(x_{0})-\F_{\inf}.
\tag{18}
\]

\item[Step 16.] \textbf{Derive the rate for the minimal mapping norm.}
Dividing \eqref{eq:18} by $K$ and using $\min_{0\le k\le K-1}\|\G_{\eta}(x_{k})\|^{2}\le \frac{1}{K}\sum_{k=0}^{K-1}\|\G_{\eta}(x_{k})\|^{2}$ yields
\[
\min_{0\le k\le K-1}\|\G_{\eta}(x_{k})\|^{2}
\le \frac{2\bigl(\F(x_{0})-\F_{\inf}\bigr)}{\eta K},
\]
which is exactly \eqref{eq:rate}. ∎

\section*{Rate Derivation (With Constants)}
The bound \eqref{eq:rate} can be rewritten as
\[
\min_{0\le k\le K-1}\|\G_{\eta}(x_{k})\|
\le \sqrt{\frac{2\bigl(\F(x_{0})-\F_{\inf}\bigr)}{\eta K}}
= O\!\bigl(K^{-1/2}\bigr).
\]
If the largest admissible stepsize $\eta=1/L$ is used, the constant becomes
\[
\sqrt{2L\bigl(\F(x_{0})-\F_{\inf}\bigr)}.
\]

\section*{Special Cases}
\begin{itemize}
  \item \textbf{Pure gradient descent ($g\equiv0$).}  
        Then $\prox_{\eta g}$ is the identity, $\G_{\eta}(x)=\grad f(x)$, and \eqref{eq:rate} reduces to the classic $O(1/\sqrt{K})$ bound on the gradient norm for smooth non‑convex $f$.
  \item \textbf{Convex $g$.}  
        The same proof holds; additionally, $\F(x_{k})$ converges to the optimal value $\F^{\star}$ (not just to a stationary point) when $f$ is also convex.
  \item \textbf{Indicator of a closed set $C$ (projection).}  
        If $g=\iota_{C}$, then $\prox_{\eta g}$ is the Euclidean projection $\Pi_{C}$. The algorithm becomes projected gradient descent, and the descent inequality \eqref{eq:16} is the standard one for non‑convex projected gradient.
\end{itemize}

\end{document}